% preamble.tex
% Shared packages and macros for all documents in this repository.
% Usage: % preamble.tex
% Shared packages and macros for all documents in this repository.
% Usage: % preamble.tex
% Shared packages and macros for all documents in this repository.
% Usage: % preamble.tex
% Shared packages and macros for all documents in this repository.
% Usage: \input{preamble} placed immediately after \documentclass.

% ══════════════════════════════════════════════════════════════════════════════
% PAGE GEOMETRY
% ══════════════════════════════════════════════════════════════════════════════

\usepackage[a4paper, margin=0.95in, top=1.1in, bottom=1.1in, headheight=14pt]{geometry}

% Professional typography
\parindent=0pt
\parskip=0.75\baselineskip plus 0.4\baselineskip minus 0.3\baselineskip
\raggedbottom

% ══════════════════════════════════════════════════════════════════════════════
% MATH PACKAGES
% ══════════════════════════════════════════════════════════════════════════════

\usepackage{amsfonts, amsmath, amssymb, amsthm}

% Optimal math spacing for readability
\thinmuskip=3mu
\medmuskip=4mu plus 1mu
\thickmuskip=5mu plus 2mu

\AtBeginDocument{
    \abovedisplayskip=1.0\baselineskip plus 0.5\baselineskip
    \abovedisplayshortskip=0.5\baselineskip plus 0.3\baselineskip
    \belowdisplayskip=1.0\baselineskip plus 0.5\baselineskip
    \belowdisplayshortskip=0.5\baselineskip plus 0.3\baselineskip
}

% ══════════════════════════════════════════════════════════════════════════════
% TYPOGRAPHY & LINE SPACING
% ══════════════════════════════════════════════════════════════════════════════

\usepackage[none]{hyphenat}
\usepackage{microtype}

% Enhanced line spacing for professional appearance
\renewcommand{\baselinestretch}{1.25}

% ══════════════════════════════════════════════════════════════════════════════
% HEADERS AND FOOTERS
% ══════════════════════════════════════════════════════════════════════════════

\usepackage{fancyhdr}
\pagestyle{fancy}
\fancyhf{}

\fancyhead[L]{\small\textsc{\leftmark}}
\fancyhead[R]{\small\textit{Discrete Mathematics Notes}}
\fancyfoot[C]{\small\thepage}

\renewcommand{\headrulewidth}{0.5pt}
\renewcommand{\footrulewidth}{0pt}

% ══════════════════════════════════════════════════════════════════════════════
% SECTION STYLING
% ══════════════════════════════════════════════════════════════════════════════

\usepackage{titlesec}

% Section: prominent, strong hierarchy
\titleformat{\section}
    {\Large\bfseries\sffamily}
    {\thesection}
    {0.7em}
    {}
    [\vspace{0.3\baselineskip}\hrule\vspace{0.4\baselineskip}]

\titlespacing{\section}
    {0pt}
    {1.8\baselineskip plus 0.5\baselineskip}
    {0.7\baselineskip plus 0.3\baselineskip}

% Subsection: clear but subordinate
\titleformat{\subsection}
    {\large\bfseries}
    {\thesubsection}
    {0.6em}
    {}

\titlespacing{\subsection}
    {0pt}
    {1.3\baselineskip plus 0.4\baselineskip}
    {0.6\baselineskip plus 0.2\baselineskip}

% Subsubsection: subtle, for fine structure
\titleformat{\subsubsection}
    {\normalsize\bfseries}
    {\thesubsubsection}
    {0.5em}
    {}

\titlespacing{\subsubsection}
    {0pt}
    {0.9\baselineskip plus 0.2\baselineskip}
    {0.4\baselineskip plus 0.1\baselineskip}

% ══════════════════════════════════════════════════════════════════════════════
% LISTS
% ══════════════════════════════════════════════════════════════════════════════

\usepackage{enumitem}
\usepackage{cancel}

\setlist[enumerate]{
    leftmargin=2.2em,
    itemsep=0.5\baselineskip,
    parsep=0.3\baselineskip,
    topsep=0.5\baselineskip
}

\setlist[itemize]{
    leftmargin=2.2em,
    itemsep=0.5\baselineskip,
    parsep=0.3\baselineskip,
    topsep=0.5\baselineskip
}

% ══════════════════════════════════════════════════════════════════════════════
% THEOREM ENVIRONMENTS
% ══════════════════════════════════════════════════════════════════════════════

\usepackage{thmtools}
\usepackage{xcolor}

% Main theorems: prominent with subtle background
\declaretheoremstyle[
    headfont=\bfseries,
    bodyfont=\normalfont,
    spaceabove=1.0\baselineskip,
    spacebelow=0.8\baselineskip,
    postheadspace=0.7em,
    headpunct=.,
    notefont=\normalfont,
    notebraces=,
]{thmstyle}

\declaretheorem[style=thmstyle, numberwithin=section]{theorem}
\declaretheorem[style=thmstyle, numberwithin=section]{proposition}
\declaretheorem[style=thmstyle, numberwithin=section]{lemma}
\declaretheorem[style=thmstyle, numberwithin=section]{corollary}

% Definitions: clearly distinguished
\declaretheoremstyle[
    headfont=\bfseries,
    bodyfont=\normalfont,
    spaceabove=0.85\baselineskip,
    spacebelow=0.7\baselineskip,
    postheadspace=0.7em,
    headpunct=.,
    notefont=\normalfont,
    notebraces=,
]{defstyle}

\declaretheorem[style=defstyle, numberwithin=section]{definition}
\declaretheorem[style=defstyle, numberwithin=section]{example}

% Remarks and notes: subtle, unnumbered
\declaretheoremstyle[
    headfont=\itshape\bfseries,
    bodyfont=\normalfont,
    spaceabove=0.6\baselineskip,
    spacebelow=0.5\baselineskip,
    postheadspace=0.6em,
    headpunct=.,
    notefont=\normalfont,
    notebraces=,
]{remstyle}

\declaretheorem[style=remstyle, unnumbered]{remark}
\declaretheorem[style=remstyle, unnumbered]{note}

% Proof environment with proper styling
\renewcommand\proofname{\bfseries Proof}
\let\oldproof\proof
\renewcommand{\proof}{%
    \oldproof[\proofname]%
    \normalfont
}

% Improved proof end symbol spacing
\renewcommand{\qedhere}{\quad\qedsymbol}

% ══════════════════════════════════════════════════════════════════════════════
% GRAPHICS & DIAGRAMS
% ══════════════════════════════════════════════════════════════════════════════

\usepackage{tikz}
\usetikzlibrary{arrows.meta, positioning, calc}

% ══════════════════════════════════════════════════════════════════════════════
% MATRICES
% ══════════════════════════════════════════════════════════════════════════════

\usepackage{blkarray}

% ══════════════════════════════════════════════════════════════════════════════
% HYPERLINKS
% ══════════════════════════════════════════════════════════════════════════════

\usepackage{hyperref}
\hypersetup{
    colorlinks=true,
    linkcolor=blue!70!black,
    urlcolor=blue!70!black,
    bookmarksnumbered=true,
    pdfborder={0 0 0}
}

% ══════════════════════════════════════════════════════════════════════════════
% NUMBER SET NOTATION
% ══════════════════════════════════════════════════════════════════════════════

\newcommand{\N}{\mathbb{N}}
\newcommand{\Z}{\mathbb{Z}}
\newcommand{\Q}{\mathbb{Q}}
\newcommand{\R}{\mathbb{R}}
\newcommand{\Ztwo}{\mathbb{Z}_2}
\newcommand{\Qstar}{\mathbb{Q}^{*}}
\newcommand{\Rplus}{\mathbb{R}^{+}}
\newcommand{\Rtwo}{\mathbb{R}^{2}}

% ══════════════════════════════════════════════════════════════════════════════
% UTILITY MACROS
% ══════════════════════════════════════════════════════════════════════════════

\newcommand{\Rel}{\mathcal{R}}
\newcommand{\defn}[1]{\textit{#1}}
 placed immediately after \documentclass.

% ══════════════════════════════════════════════════════════════════════════════
% PAGE GEOMETRY
% ══════════════════════════════════════════════════════════════════════════════

\usepackage[a4paper, margin=0.95in, top=1.1in, bottom=1.1in, headheight=14pt]{geometry}

% Professional typography
\parindent=0pt
\parskip=0.75\baselineskip plus 0.4\baselineskip minus 0.3\baselineskip
\raggedbottom

% ══════════════════════════════════════════════════════════════════════════════
% MATH PACKAGES
% ══════════════════════════════════════════════════════════════════════════════

\usepackage{amsfonts, amsmath, amssymb, amsthm}

% Optimal math spacing for readability
\thinmuskip=3mu
\medmuskip=4mu plus 1mu
\thickmuskip=5mu plus 2mu

\AtBeginDocument{
    \abovedisplayskip=1.0\baselineskip plus 0.5\baselineskip
    \abovedisplayshortskip=0.5\baselineskip plus 0.3\baselineskip
    \belowdisplayskip=1.0\baselineskip plus 0.5\baselineskip
    \belowdisplayshortskip=0.5\baselineskip plus 0.3\baselineskip
}

% ══════════════════════════════════════════════════════════════════════════════
% TYPOGRAPHY & LINE SPACING
% ══════════════════════════════════════════════════════════════════════════════

\usepackage[none]{hyphenat}
\usepackage{microtype}

% Enhanced line spacing for professional appearance
\renewcommand{\baselinestretch}{1.25}

% ══════════════════════════════════════════════════════════════════════════════
% HEADERS AND FOOTERS
% ══════════════════════════════════════════════════════════════════════════════

\usepackage{fancyhdr}
\pagestyle{fancy}
\fancyhf{}

\fancyhead[L]{\small\textsc{\leftmark}}
\fancyhead[R]{\small\textit{Discrete Mathematics Notes}}
\fancyfoot[C]{\small\thepage}

\renewcommand{\headrulewidth}{0.5pt}
\renewcommand{\footrulewidth}{0pt}

% ══════════════════════════════════════════════════════════════════════════════
% SECTION STYLING
% ══════════════════════════════════════════════════════════════════════════════

\usepackage{titlesec}

% Section: prominent, strong hierarchy
\titleformat{\section}
    {\Large\bfseries\sffamily}
    {\thesection}
    {0.7em}
    {}
    [\vspace{0.3\baselineskip}\hrule\vspace{0.4\baselineskip}]

\titlespacing{\section}
    {0pt}
    {1.8\baselineskip plus 0.5\baselineskip}
    {0.7\baselineskip plus 0.3\baselineskip}

% Subsection: clear but subordinate
\titleformat{\subsection}
    {\large\bfseries}
    {\thesubsection}
    {0.6em}
    {}

\titlespacing{\subsection}
    {0pt}
    {1.3\baselineskip plus 0.4\baselineskip}
    {0.6\baselineskip plus 0.2\baselineskip}

% Subsubsection: subtle, for fine structure
\titleformat{\subsubsection}
    {\normalsize\bfseries}
    {\thesubsubsection}
    {0.5em}
    {}

\titlespacing{\subsubsection}
    {0pt}
    {0.9\baselineskip plus 0.2\baselineskip}
    {0.4\baselineskip plus 0.1\baselineskip}

% ══════════════════════════════════════════════════════════════════════════════
% LISTS
% ══════════════════════════════════════════════════════════════════════════════

\usepackage{enumitem}
\usepackage{cancel}

\setlist[enumerate]{
    leftmargin=2.2em,
    itemsep=0.5\baselineskip,
    parsep=0.3\baselineskip,
    topsep=0.5\baselineskip
}

\setlist[itemize]{
    leftmargin=2.2em,
    itemsep=0.5\baselineskip,
    parsep=0.3\baselineskip,
    topsep=0.5\baselineskip
}

% ══════════════════════════════════════════════════════════════════════════════
% THEOREM ENVIRONMENTS
% ══════════════════════════════════════════════════════════════════════════════

\usepackage{thmtools}
\usepackage{xcolor}

% Main theorems: prominent with subtle background
\declaretheoremstyle[
    headfont=\bfseries,
    bodyfont=\normalfont,
    spaceabove=1.0\baselineskip,
    spacebelow=0.8\baselineskip,
    postheadspace=0.7em,
    headpunct=.,
    notefont=\normalfont,
    notebraces=,
]{thmstyle}

\declaretheorem[style=thmstyle, numberwithin=section]{theorem}
\declaretheorem[style=thmstyle, numberwithin=section]{proposition}
\declaretheorem[style=thmstyle, numberwithin=section]{lemma}
\declaretheorem[style=thmstyle, numberwithin=section]{corollary}

% Definitions: clearly distinguished
\declaretheoremstyle[
    headfont=\bfseries,
    bodyfont=\normalfont,
    spaceabove=0.85\baselineskip,
    spacebelow=0.7\baselineskip,
    postheadspace=0.7em,
    headpunct=.,
    notefont=\normalfont,
    notebraces=,
]{defstyle}

\declaretheorem[style=defstyle, numberwithin=section]{definition}
\declaretheorem[style=defstyle, numberwithin=section]{example}

% Remarks and notes: subtle, unnumbered
\declaretheoremstyle[
    headfont=\itshape\bfseries,
    bodyfont=\normalfont,
    spaceabove=0.6\baselineskip,
    spacebelow=0.5\baselineskip,
    postheadspace=0.6em,
    headpunct=.,
    notefont=\normalfont,
    notebraces=,
]{remstyle}

\declaretheorem[style=remstyle, unnumbered]{remark}
\declaretheorem[style=remstyle, unnumbered]{note}

% Proof environment with proper styling
\renewcommand\proofname{\bfseries Proof}
\let\oldproof\proof
\renewcommand{\proof}{%
    \oldproof[\proofname]%
    \normalfont
}

% Improved proof end symbol spacing
\renewcommand{\qedhere}{\quad\qedsymbol}

% ══════════════════════════════════════════════════════════════════════════════
% GRAPHICS & DIAGRAMS
% ══════════════════════════════════════════════════════════════════════════════

\usepackage{tikz}
\usetikzlibrary{arrows.meta, positioning, calc}

% ══════════════════════════════════════════════════════════════════════════════
% MATRICES
% ══════════════════════════════════════════════════════════════════════════════

\usepackage{blkarray}

% ══════════════════════════════════════════════════════════════════════════════
% HYPERLINKS
% ══════════════════════════════════════════════════════════════════════════════

\usepackage{hyperref}
\hypersetup{
    colorlinks=true,
    linkcolor=blue!70!black,
    urlcolor=blue!70!black,
    bookmarksnumbered=true,
    pdfborder={0 0 0}
}

% ══════════════════════════════════════════════════════════════════════════════
% NUMBER SET NOTATION
% ══════════════════════════════════════════════════════════════════════════════

\newcommand{\N}{\mathbb{N}}
\newcommand{\Z}{\mathbb{Z}}
\newcommand{\Q}{\mathbb{Q}}
\newcommand{\R}{\mathbb{R}}
\newcommand{\Ztwo}{\mathbb{Z}_2}
\newcommand{\Qstar}{\mathbb{Q}^{*}}
\newcommand{\Rplus}{\mathbb{R}^{+}}
\newcommand{\Rtwo}{\mathbb{R}^{2}}

% ══════════════════════════════════════════════════════════════════════════════
% UTILITY MACROS
% ══════════════════════════════════════════════════════════════════════════════

\newcommand{\Rel}{\mathcal{R}}
\newcommand{\defn}[1]{\textit{#1}}
 placed immediately after \documentclass.

% ══════════════════════════════════════════════════════════════════════════════
% PAGE GEOMETRY
% ══════════════════════════════════════════════════════════════════════════════

\usepackage[a4paper, margin=0.95in, top=1.1in, bottom=1.1in, headheight=14pt]{geometry}

% Professional typography
\parindent=0pt
\parskip=0.75\baselineskip plus 0.4\baselineskip minus 0.3\baselineskip
\raggedbottom

% ══════════════════════════════════════════════════════════════════════════════
% MATH PACKAGES
% ══════════════════════════════════════════════════════════════════════════════

\usepackage{amsfonts, amsmath, amssymb, amsthm}

% Optimal math spacing for readability
\thinmuskip=3mu
\medmuskip=4mu plus 1mu
\thickmuskip=5mu plus 2mu

\AtBeginDocument{
    \abovedisplayskip=1.0\baselineskip plus 0.5\baselineskip
    \abovedisplayshortskip=0.5\baselineskip plus 0.3\baselineskip
    \belowdisplayskip=1.0\baselineskip plus 0.5\baselineskip
    \belowdisplayshortskip=0.5\baselineskip plus 0.3\baselineskip
}

% ══════════════════════════════════════════════════════════════════════════════
% TYPOGRAPHY & LINE SPACING
% ══════════════════════════════════════════════════════════════════════════════

\usepackage[none]{hyphenat}
\usepackage{microtype}

% Enhanced line spacing for professional appearance
\renewcommand{\baselinestretch}{1.25}

% ══════════════════════════════════════════════════════════════════════════════
% HEADERS AND FOOTERS
% ══════════════════════════════════════════════════════════════════════════════

\usepackage{fancyhdr}
\pagestyle{fancy}
\fancyhf{}

\fancyhead[L]{\small\textsc{\leftmark}}
\fancyhead[R]{\small\textit{Discrete Mathematics Notes}}
\fancyfoot[C]{\small\thepage}

\renewcommand{\headrulewidth}{0.5pt}
\renewcommand{\footrulewidth}{0pt}

% ══════════════════════════════════════════════════════════════════════════════
% SECTION STYLING
% ══════════════════════════════════════════════════════════════════════════════

\usepackage{titlesec}

% Section: prominent, strong hierarchy
\titleformat{\section}
    {\Large\bfseries\sffamily}
    {\thesection}
    {0.7em}
    {}
    [\vspace{0.3\baselineskip}\hrule\vspace{0.4\baselineskip}]

\titlespacing{\section}
    {0pt}
    {1.8\baselineskip plus 0.5\baselineskip}
    {0.7\baselineskip plus 0.3\baselineskip}

% Subsection: clear but subordinate
\titleformat{\subsection}
    {\large\bfseries}
    {\thesubsection}
    {0.6em}
    {}

\titlespacing{\subsection}
    {0pt}
    {1.3\baselineskip plus 0.4\baselineskip}
    {0.6\baselineskip plus 0.2\baselineskip}

% Subsubsection: subtle, for fine structure
\titleformat{\subsubsection}
    {\normalsize\bfseries}
    {\thesubsubsection}
    {0.5em}
    {}

\titlespacing{\subsubsection}
    {0pt}
    {0.9\baselineskip plus 0.2\baselineskip}
    {0.4\baselineskip plus 0.1\baselineskip}

% ══════════════════════════════════════════════════════════════════════════════
% LISTS
% ══════════════════════════════════════════════════════════════════════════════

\usepackage{enumitem}
\usepackage{cancel}

\setlist[enumerate]{
    leftmargin=2.2em,
    itemsep=0.5\baselineskip,
    parsep=0.3\baselineskip,
    topsep=0.5\baselineskip
}

\setlist[itemize]{
    leftmargin=2.2em,
    itemsep=0.5\baselineskip,
    parsep=0.3\baselineskip,
    topsep=0.5\baselineskip
}

% ══════════════════════════════════════════════════════════════════════════════
% THEOREM ENVIRONMENTS
% ══════════════════════════════════════════════════════════════════════════════

\usepackage{thmtools}
\usepackage{xcolor}

% Main theorems: prominent with subtle background
\declaretheoremstyle[
    headfont=\bfseries,
    bodyfont=\normalfont,
    spaceabove=1.0\baselineskip,
    spacebelow=0.8\baselineskip,
    postheadspace=0.7em,
    headpunct=.,
    notefont=\normalfont,
    notebraces=,
]{thmstyle}

\declaretheorem[style=thmstyle, numberwithin=section]{theorem}
\declaretheorem[style=thmstyle, numberwithin=section]{proposition}
\declaretheorem[style=thmstyle, numberwithin=section]{lemma}
\declaretheorem[style=thmstyle, numberwithin=section]{corollary}

% Definitions: clearly distinguished
\declaretheoremstyle[
    headfont=\bfseries,
    bodyfont=\normalfont,
    spaceabove=0.85\baselineskip,
    spacebelow=0.7\baselineskip,
    postheadspace=0.7em,
    headpunct=.,
    notefont=\normalfont,
    notebraces=,
]{defstyle}

\declaretheorem[style=defstyle, numberwithin=section]{definition}
\declaretheorem[style=defstyle, numberwithin=section]{example}

% Remarks and notes: subtle, unnumbered
\declaretheoremstyle[
    headfont=\itshape\bfseries,
    bodyfont=\normalfont,
    spaceabove=0.6\baselineskip,
    spacebelow=0.5\baselineskip,
    postheadspace=0.6em,
    headpunct=.,
    notefont=\normalfont,
    notebraces=,
]{remstyle}

\declaretheorem[style=remstyle, unnumbered]{remark}
\declaretheorem[style=remstyle, unnumbered]{note}

% Proof environment with proper styling
\renewcommand\proofname{\bfseries Proof}
\let\oldproof\proof
\renewcommand{\proof}{%
    \oldproof[\proofname]%
    \normalfont
}

% Improved proof end symbol spacing
\renewcommand{\qedhere}{\quad\qedsymbol}

% ══════════════════════════════════════════════════════════════════════════════
% GRAPHICS & DIAGRAMS
% ══════════════════════════════════════════════════════════════════════════════

\usepackage{tikz}
\usetikzlibrary{arrows.meta, positioning, calc}

% ══════════════════════════════════════════════════════════════════════════════
% MATRICES
% ══════════════════════════════════════════════════════════════════════════════

\usepackage{blkarray}

% ══════════════════════════════════════════════════════════════════════════════
% HYPERLINKS
% ══════════════════════════════════════════════════════════════════════════════

\usepackage{hyperref}
\hypersetup{
    colorlinks=true,
    linkcolor=blue!70!black,
    urlcolor=blue!70!black,
    bookmarksnumbered=true,
    pdfborder={0 0 0}
}

% ══════════════════════════════════════════════════════════════════════════════
% NUMBER SET NOTATION
% ══════════════════════════════════════════════════════════════════════════════

\newcommand{\N}{\mathbb{N}}
\newcommand{\Z}{\mathbb{Z}}
\newcommand{\Q}{\mathbb{Q}}
\newcommand{\R}{\mathbb{R}}
\newcommand{\Ztwo}{\mathbb{Z}_2}
\newcommand{\Qstar}{\mathbb{Q}^{*}}
\newcommand{\Rplus}{\mathbb{R}^{+}}
\newcommand{\Rtwo}{\mathbb{R}^{2}}

% ══════════════════════════════════════════════════════════════════════════════
% UTILITY MACROS
% ══════════════════════════════════════════════════════════════════════════════

\newcommand{\Rel}{\mathcal{R}}
\newcommand{\defn}[1]{\textit{#1}}
 placed immediately after \documentclass.

% ══════════════════════════════════════════════════════════════════════════════
% PAGE GEOMETRY
% ══════════════════════════════════════════════════════════════════════════════

\usepackage[a4paper, margin=0.95in, top=1.1in, bottom=1.1in, headheight=14pt]{geometry}

% Professional typography
\parindent=0pt
\parskip=0.75\baselineskip plus 0.4\baselineskip minus 0.3\baselineskip
\raggedbottom

% ══════════════════════════════════════════════════════════════════════════════
% MATH PACKAGES
% ══════════════════════════════════════════════════════════════════════════════

\usepackage{amsfonts, amsmath, amssymb, amsthm}

% Optimal math spacing for readability
\thinmuskip=3mu
\medmuskip=4mu plus 1mu
\thickmuskip=5mu plus 2mu

\AtBeginDocument{
    \abovedisplayskip=1.0\baselineskip plus 0.5\baselineskip
    \abovedisplayshortskip=0.5\baselineskip plus 0.3\baselineskip
    \belowdisplayskip=1.0\baselineskip plus 0.5\baselineskip
    \belowdisplayshortskip=0.5\baselineskip plus 0.3\baselineskip
}

% ══════════════════════════════════════════════════════════════════════════════
% TYPOGRAPHY & LINE SPACING
% ══════════════════════════════════════════════════════════════════════════════

\usepackage[none]{hyphenat}
\usepackage{microtype}

% Enhanced line spacing for professional appearance
\renewcommand{\baselinestretch}{1.25}

% ══════════════════════════════════════════════════════════════════════════════
% HEADERS AND FOOTERS
% ══════════════════════════════════════════════════════════════════════════════

\usepackage{fancyhdr}
\pagestyle{fancy}
\fancyhf{}

\fancyhead[L]{\small\textsc{\leftmark}}
\fancyhead[R]{\small\textit{Discrete Mathematics Notes}}
\fancyfoot[C]{\small\thepage}

\renewcommand{\headrulewidth}{0.5pt}
\renewcommand{\footrulewidth}{0pt}

% ══════════════════════════════════════════════════════════════════════════════
% SECTION STYLING
% ══════════════════════════════════════════════════════════════════════════════

\usepackage{titlesec}

% Section: prominent, strong hierarchy
\titleformat{\section}
    {\Large\bfseries\sffamily}
    {\thesection}
    {0.7em}
    {}
    [\vspace{0.3\baselineskip}\hrule\vspace{0.4\baselineskip}]

\titlespacing{\section}
    {0pt}
    {1.8\baselineskip plus 0.5\baselineskip}
    {0.7\baselineskip plus 0.3\baselineskip}

% Subsection: clear but subordinate
\titleformat{\subsection}
    {\large\bfseries}
    {\thesubsection}
    {0.6em}
    {}

\titlespacing{\subsection}
    {0pt}
    {1.3\baselineskip plus 0.4\baselineskip}
    {0.6\baselineskip plus 0.2\baselineskip}

% Subsubsection: subtle, for fine structure
\titleformat{\subsubsection}
    {\normalsize\bfseries}
    {\thesubsubsection}
    {0.5em}
    {}

\titlespacing{\subsubsection}
    {0pt}
    {0.9\baselineskip plus 0.2\baselineskip}
    {0.4\baselineskip plus 0.1\baselineskip}

% ══════════════════════════════════════════════════════════════════════════════
% LISTS
% ══════════════════════════════════════════════════════════════════════════════

\usepackage{enumitem}
\usepackage{cancel}

\setlist[enumerate]{
    leftmargin=2.2em,
    itemsep=0.5\baselineskip,
    parsep=0.3\baselineskip,
    topsep=0.5\baselineskip
}

\setlist[itemize]{
    leftmargin=2.2em,
    itemsep=0.5\baselineskip,
    parsep=0.3\baselineskip,
    topsep=0.5\baselineskip
}

% ══════════════════════════════════════════════════════════════════════════════
% THEOREM ENVIRONMENTS
% ══════════════════════════════════════════════════════════════════════════════

\usepackage{thmtools}
\usepackage{xcolor}

% Main theorems: prominent with subtle background
\declaretheoremstyle[
    headfont=\bfseries,
    bodyfont=\normalfont,
    spaceabove=1.0\baselineskip,
    spacebelow=0.8\baselineskip,
    postheadspace=0.7em,
    headpunct=.,
    notefont=\normalfont,
    notebraces=,
]{thmstyle}

\declaretheorem[style=thmstyle, numberwithin=section]{theorem}
\declaretheorem[style=thmstyle, numberwithin=section]{proposition}
\declaretheorem[style=thmstyle, numberwithin=section]{lemma}
\declaretheorem[style=thmstyle, numberwithin=section]{corollary}

% Definitions: clearly distinguished
\declaretheoremstyle[
    headfont=\bfseries,
    bodyfont=\normalfont,
    spaceabove=0.85\baselineskip,
    spacebelow=0.7\baselineskip,
    postheadspace=0.7em,
    headpunct=.,
    notefont=\normalfont,
    notebraces=,
]{defstyle}

\declaretheorem[style=defstyle, numberwithin=section]{definition}
\declaretheorem[style=defstyle, numberwithin=section]{example}

% Remarks and notes: subtle, unnumbered
\declaretheoremstyle[
    headfont=\itshape\bfseries,
    bodyfont=\normalfont,
    spaceabove=0.6\baselineskip,
    spacebelow=0.5\baselineskip,
    postheadspace=0.6em,
    headpunct=.,
    notefont=\normalfont,
    notebraces=,
]{remstyle}

\declaretheorem[style=remstyle, unnumbered]{remark}
\declaretheorem[style=remstyle, unnumbered]{note}

% Proof environment with proper styling
\renewcommand\proofname{\bfseries Proof}
\let\oldproof\proof
\renewcommand{\proof}{%
    \oldproof[\proofname]%
    \normalfont
}

% Improved proof end symbol spacing
\renewcommand{\qedhere}{\quad\qedsymbol}

% ══════════════════════════════════════════════════════════════════════════════
% GRAPHICS & DIAGRAMS
% ══════════════════════════════════════════════════════════════════════════════

\usepackage{tikz}
\usetikzlibrary{arrows.meta, positioning, calc}

% ══════════════════════════════════════════════════════════════════════════════
% MATRICES
% ══════════════════════════════════════════════════════════════════════════════

\usepackage{blkarray}

% ══════════════════════════════════════════════════════════════════════════════
% HYPERLINKS
% ══════════════════════════════════════════════════════════════════════════════

\usepackage{hyperref}
\hypersetup{
    colorlinks=true,
    linkcolor=blue!70!black,
    urlcolor=blue!70!black,
    bookmarksnumbered=true,
    pdfborder={0 0 0}
}

% ══════════════════════════════════════════════════════════════════════════════
% NUMBER SET NOTATION
% ══════════════════════════════════════════════════════════════════════════════

\newcommand{\N}{\mathbb{N}}
\newcommand{\Z}{\mathbb{Z}}
\newcommand{\Q}{\mathbb{Q}}
\newcommand{\R}{\mathbb{R}}
\newcommand{\Ztwo}{\mathbb{Z}_2}
\newcommand{\Qstar}{\mathbb{Q}^{*}}
\newcommand{\Rplus}{\mathbb{R}^{+}}
\newcommand{\Rtwo}{\mathbb{R}^{2}}

% ══════════════════════════════════════════════════════════════════════════════
% UTILITY MACROS
% ══════════════════════════════════════════════════════════════════════════════

\newcommand{\Rel}{\mathcal{R}}
\newcommand{\defn}[1]{\textit{#1}}
